\documentclass[aspectratio=169,11pt]{beamer}
\newcommand{\footerlabel}{DEMO102 Networks \textbullet{} Lecture 2 (slides continue from Lecture 1)}
% 共用样式 —— 虚构示例课件（DEMO），内容全部自编，不来自任何真实课程。
% 用法：每份 deck 先 \documentclass，再定义 \footerlabel，最后 % 共用样式 —— 虚构示例课件（DEMO），内容全部自编，不来自任何真实课程。
% 用法：每份 deck 先 \documentclass，再定义 \footerlabel，最后 % 共用样式 —— 虚构示例课件（DEMO），内容全部自编，不来自任何真实课程。
% 用法：每份 deck 先 \documentclass，再定义 \footerlabel，最后 \input{_style}。
\usepackage[T1]{fontenc}
\usepackage{lmodern}
\usepackage{amsmath,amssymb}
\usetheme{default}
\setbeamertemplate{navigation symbols}{}
\setbeamercolor{frametitle}{fg=blue!50!black}
\setbeamerfont{frametitle}{series=\bfseries,size=\Large}
\setbeamercolor{title}{fg=blue!50!black}
\setbeamercolor{footline}{fg=black!60}
\setbeamertemplate{footline}{%
  \leavevmode\hbox{%
    \begin{beamercolorbox}[wd=.85\paperwidth,ht=2.5ex,dp=1ex,left,leftskip=1em]{footline}\footnotesize\footerlabel\end{beamercolorbox}%
    \begin{beamercolorbox}[wd=.15\paperwidth,ht=2.5ex,dp=1ex,right,rightskip=1em]{footline}\footnotesize\insertframenumber\end{beamercolorbox}}%
  \vskip0pt}
\setbeamertemplate{itemize items}[circle]
。
\usepackage[T1]{fontenc}
\usepackage{lmodern}
\usepackage{amsmath,amssymb}
\usetheme{default}
\setbeamertemplate{navigation symbols}{}
\setbeamercolor{frametitle}{fg=blue!50!black}
\setbeamerfont{frametitle}{series=\bfseries,size=\Large}
\setbeamercolor{title}{fg=blue!50!black}
\setbeamercolor{footline}{fg=black!60}
\setbeamertemplate{footline}{%
  \leavevmode\hbox{%
    \begin{beamercolorbox}[wd=.85\paperwidth,ht=2.5ex,dp=1ex,left,leftskip=1em]{footline}\footnotesize\footerlabel\end{beamercolorbox}%
    \begin{beamercolorbox}[wd=.15\paperwidth,ht=2.5ex,dp=1ex,right,rightskip=1em]{footline}\footnotesize\insertframenumber\end{beamercolorbox}}%
  \vskip0pt}
\setbeamertemplate{itemize items}[circle]
。
\usepackage[T1]{fontenc}
\usepackage{lmodern}
\usepackage{amsmath,amssymb}
\usetheme{default}
\setbeamertemplate{navigation symbols}{}
\setbeamercolor{frametitle}{fg=blue!50!black}
\setbeamerfont{frametitle}{series=\bfseries,size=\Large}
\setbeamercolor{title}{fg=blue!50!black}
\setbeamercolor{footline}{fg=black!60}
\setbeamertemplate{footline}{%
  \leavevmode\hbox{%
    \begin{beamercolorbox}[wd=.85\paperwidth,ht=2.5ex,dp=1ex,left,leftskip=1em]{footline}\footnotesize\footerlabel\end{beamercolorbox}%
    \begin{beamercolorbox}[wd=.15\paperwidth,ht=2.5ex,dp=1ex,right,rightskip=1em]{footline}\footnotesize\insertframenumber\end{beamercolorbox}}%
  \vskip0pt}
\setbeamertemplate{itemize items}[circle]

\title{Computer Networks}
\subtitle{Lecture 2: The Physical Layer}
\author{Demo Lecturer}
\date{Semester 2, 2026}

\begin{document}

\begin{frame}
  \titlepage
\end{frame}
% 续篇：本讲的页脚号接着 Lecture 1 的 7 页往下编（示例：封面印 1，正文从 8 起 —— 判续篇要看 p.2 不看 p.1）
\setcounter{framenumber}{7}

\begin{frame}{Recap: Where We Are}
  \begin{itemize}
    \item Five layers: application, transport, network, link, \textbf{physical}
    \item The physical layer moves raw \textbf{bits} across a medium
    \item Today: how fast can a link carry bits?
  \end{itemize}
\end{frame}

\begin{frame}{Signals and Bandwidth}
  \begin{itemize}
    \item A signal varies over time; its content spans a range of frequencies
    \item \textbf{Bandwidth} $B$ = the width of the usable frequency range, in hertz (Hz)
    \item A channel passes frequencies from $f_1$ to $f_2$, so $B = f_2 - f_1$
  \end{itemize}
  \medskip
  Example: a telephone line passing 300 Hz to 3300 Hz has $B = 3000$ Hz.
\end{frame}

\begin{frame}{Bandwidth, Data Rate, Latency}
  \begin{center}
  \begin{tabular}{lll}
    \hline
    Quantity & Unit & Meaning \\
    \hline
    Bandwidth $B$ & Hz & width of the frequency range \\
    Data rate & bits/s & how many bits per second we send \\
    Latency & s & time for one bit to get across \\
    \hline
  \end{tabular}
  \end{center}
  \medskip
  Careful: in everyday speech ``bandwidth'' is often used for the data rate. In this subject
  they are different quantities.
\end{frame}

\begin{frame}{Nyquist: A Noiseless Channel}
  For a noiseless channel of bandwidth $B$ Hz and signals with $V$ distinct levels:
  \[
    \text{max data rate} = 2B\log_2 V \text{ bits/s}
  \]
  \textbf{Example.} $B = 3000$ Hz, binary signal ($V = 2$):
  \[
    2 \times 3000 \times \log_2 2 = 6000 \text{ bits/s}
  \]
  More signal levels give more bits per symbol -- but noise limits how many levels we can tell apart.
\end{frame}

\begin{frame}{Shannon: A Noisy Channel}
  With signal-to-noise ratio $S/N$ (a plain ratio, not decibels):
  \[
    C = B\log_2\!\left(1 + \frac{S}{N}\right) \text{ bits/s}
  \]
  \textbf{Example.} $B = 3000$ Hz and $30$ dB: $S/N = 10^{30/10} = 1000$, so
  \[
    C = 3000 \times \log_2(1001) \approx 29\,900 \text{ bits/s}
  \]
  No coding scheme can beat $C$, however many signal levels are used.
\end{frame}

\begin{frame}{Transmission Media}
  \begin{center}
  \begin{tabular}{lll}
    \hline
    Medium & Typical use & Main limit \\
    \hline
    Twisted pair & office LAN, last mile & attenuation over distance \\
    Optical fibre & backbone links & cost of installation \\
    Radio & Wi-Fi, mobile & shared spectrum, interference \\
    \hline
  \end{tabular}
  \end{center}
\end{frame}

\begin{frame}{Summary}
  \begin{itemize}
    \item Bandwidth ($B$, Hz) and data rate (bits/s) are different quantities
    \item Nyquist: $2B\log_2 V$ for a noiseless channel
    \item Shannon: $B\log_2(1 + S/N)$ for a noisy channel -- an upper limit
    \item Next: the link layer and how a shared medium is accessed
  \end{itemize}
\end{frame}

\end{document}
